\documentclass[10pt,a4paper]{article}
\usepackage[utf8]{inputenc}
\usepackage[T1]{fontenc}
\IfFileExists{brazil.ldf}{\usepackage[brazil]{babel}}{}
\IfFileExists{charter.sty}{\usepackage{charter}}{\usepackage{lmodern}}
\IfFileExists{charter.sty}{\renewcommand{\ttdefault}{pcr}}{}
\usepackage[a4paper,top=1.1cm,bottom=1.3cm,left=1.5cm,right=1.5cm,footskip=0.65cm]{geometry}
\usepackage{microtype}
\usepackage{amsmath}
\usepackage{booktabs}
\usepackage{tabularx}
\usepackage{enumitem}
\usepackage[dvipsnames]{xcolor}
\usepackage{titlesec}
\usepackage{fancyhdr}
\usepackage[hidelinks]{hyperref}
\usepackage{pgfplots}
\pgfplotsset{compat=1.18}
\usepackage{xstring}
\usepackage[font=small,labelfont=bf,labelsep=colon,skip=2pt]{caption}
\captionsetup[figure]{name=Figura}
\setlength{\abovedisplayskip}{2pt}\setlength{\belowdisplayskip}{2pt}
\setlength{\abovedisplayshortskip}{1pt}\setlength{\belowdisplayshortskip}{1pt}
\setlength{\textfloatsep}{5pt}\setlength{\floatsep}{5pt}\setlength{\intextsep}{4pt}

% Gerado por scripts/gerar_numeros_latex_2t.py. NAO editar a mao.
\newcommand{\horacorte}{sábado, 24/10/2026, às 20h00}
\newcommand{\statusprevisao}{preliminar}
\newcommand{\pFlavio}{51,3}
\newcommand{\pLula}{48,7}
\newcommand{\pAbst}{21,1}
\newcommand{\pBrancos}{1,8}
\newcommand{\pNulos}{2,9}
\newcommand{\nPesq}{26}
\newcommand{\nInst}{14}
\newcommand{\nInstUmT}{5}
\newcommand{\nHist}{111}
\newcommand{\dataUltimaPesq}{21/10}
\newcommand{\sigmaZero}{1,5}
\newcommand{\medErroUmT}{2,42}
\newcommand{\maeVenc}{0,65}
\newcommand{\maeMZero}{0,92}
\newcommand{\deltaVenc}{0,27}
\newcommand{\seVenc}{0,14}
\newcommand{\loeoVenc}{0,48}
\newcommand{\loeoMZero}{0,77}
\newcommand{\lulaUmT}{45,16}
\newcommand{\flavioUmT}{47,03}
\newcommand{\tercUmT}{7,81}
\newcommand{\flavioDoisCand}{51,0}
\newcommand{\lulaDoisCand}{49,0}
\newcommand{\tauMediana}{0,45}
\newcommand{\tauMedia}{0,61}
\newcommand{\tauImplFlavio}{55}
\newcommand{\tauImplLula}{45}
\newcommand{\priorMedianaFlavio}{51,4}
\newcommand{\priorMediaFlavio}{50,1}
\newcommand{\sensMZero}{51,3}
\newcommand{\sensMUm}{51,1}
\newcommand{\sensPersist}{51,7}
\newcommand{\sensMin}{50,1}
\newcommand{\sensMax}{51,8}
\newcommand{\nSens}{11}
\newcommand{\abstUmT}{21,08}
\newcommand{\brancosUmT}{1,84}
\newcommand{\nulosUmT}{2,93}
\newcommand{\govShareVinteSeis}{7,9}
\newcommand{\govNVinteSeis}{6}
\newcommand{\govShareHist}{41,6}
\newcommand{\deltaAbstCom}{1,72}
\newcommand{\deltaAbstSem}{1,60}
\newcommand{\pesoMax}{13,4}
\newcommand{\instPesoMax}{Quaest}
\newcommand{\pesoMin}{3,5}
\newcommand{\instPesoMin}{CNT/MDA}
\newcommand{\rotulosPesos}{{CNT/MDA},{Real Time Big Data},{Datafolha},{Palver},{FSB/BTG},{Alfa Inteligência},{Ideia},{Futura},{Gerp},{PoderData},{Veritá},{Ipespe},{AtlasIntel},{Quaest}}
\newcommand{\nPesos}{13}


\definecolor{azul}{HTML}{1F3A5F}
\definecolor{cinza}{HTML}{555555}
\definecolor{vermelho}{HTML}{B03A2E}

\titleformat{\section}{\normalfont\bfseries\color{azul}}{\thesection.}{0.5em}{}
\titlespacing*{\section}{0pt}{3pt}{1pt}
\setlist{nosep,leftmargin=1.2em,topsep=1pt}
\setlength{\parindent}{0pt}
\setlength{\parskip}{2pt}
\renewcommand{\arraystretch}{1.0}
\setlength{\tabcolsep}{3.5pt}

\pagestyle{fancy}
\fancyhf{}
\renewcommand{\headrulewidth}{0pt}
\fancyfoot[L]{\footnotesize\color{cinza}Desafio de Estatística e Econometria, FGV EPGE}
\fancyfoot[R]{\footnotesize\color{cinza}\thepage/2}

\StrSubstitute{\pFlavio}{,}{.}[\pFlavioDot]
\StrSubstitute{\flavioDoisCand}{,}{.}[\flavioDoisCandDot]

\begin{document}

{\large\bfseries\color{azul} Previsão do 2º turno presidencial de 2026: nota metodológica}\\[1pt]
{\small João Pedro Valuche de Andrade Pereira \textbullet\ Lethicia Manfioletti Possamai \textbullet\ Arthur Caron Lyra}\\[2pt]
\colorbox{azul!8}{\parbox{\dimexpr\linewidth-2\fboxsep}{\small
\textbf{Horário de corte dos dados: \horacorte{} (horário de Brasília).} Nenhuma pesquisa divulgada depois desse horário entrou na previsão. Status destes números: \statusprevisao{} (\nPesq{} pesquisas, a última divulgada em \dataUltimaPesq).}}

\vspace{2pt}
{\setlength{\tabcolsep}{0pt}\small
\begin{tabularx}{\linewidth}{@{}*{5}{>{\centering\arraybackslash}X}@{}}
\textcolor{azul}{\Large\bfseries \pFlavio\%} & \textcolor{azul}{\Large\bfseries \pLula\%} & \textcolor{cinza}{\Large\bfseries \pAbst\%} & \textcolor{cinza}{\Large\bfseries \pBrancos\%} & \textcolor{cinza}{\Large\bfseries \pNulos\%}\\
\footnotesize Flávio Bolsonaro (PL) & \footnotesize Lula (PT) & \footnotesize Abstenção & \footnotesize Brancos & \footnotesize Nulos\\[-1pt]
\multicolumn{2}{c}{\scriptsize\color{cinza}\% dos votos válidos} & \multicolumn{3}{c}{\scriptsize\color{cinza}\% do eleitorado apto (abstenção) e do comparecimento (brancos e nulos)}
\end{tabularx}}

\section{Objetivo e métrica}
Com dois candidatos, os votos válidos somam 100 e o erro de um é o erro do outro com sinal trocado. O MAE da previsão é, portanto, o erro absoluto de um candidato, que é também metade do erro na margem:
$\text{MAE}=\tfrac12\sum_{k=1}^{2}|\hat V_k-V_k|=|\hat V_{\text{Lula}}-V_{\text{Lula}}|$.
Errar a margem em 2 pontos custa 1 ponto de MAE. Toda a modelagem se reduz a estimar bem uma única proporção.

\section{Dados}
\textbf{1º turno de 2026.} Resultado oficial do TSE por município e zona (arquivos do portal de dados abertos, com hash SHA-256): Flávio \flavioUmT\%, Lula \lulaUmT\%, demais candidatos \tercUmT\% dos válidos; abstenção \abstUmT\%, brancos \brancosUmT\% e nulos \nulosUmT\%.
\textbf{Histórico.} Resultados de 1º e 2º turno de 2006 a 2022, por UF, e \nHist{} pesquisas nacionais de 2º turno divulgadas entre os dois turnos dessas eleições.
\textbf{Pesquisas de 2026.} \nPesq{} pesquisas estimuladas de 2º turno de \nInst{} institutos, todas com registro no TSE e fonte anotada na base, divulgadas a partir de 05/10. Convertemos cada uma em votos válidos, $\hat v_{j}=P_{j,\text{Lula}}/(P_{j,\text{Lula}}+P_{j,\text{Flávio}})$, o que exclui brancos e nulos e distribui os indecisos na proporção dos decididos. De cada instituto usamos só a pesquisa mais recente, para que quem divulga com mais frequência não pese mais.

\section{Modelos e validação}
Testamos três agregadores. O M0 é a média simples das últimas pesquisas de cada instituto. O M1 usa todas as pesquisas da janela com peso $w_i=\sqrt{N_i}\,2^{-\Delta t_i/h}$, que favorece amostras maiores e pesquisas recentes ($h\in\{3,5,7\}$ dias). O M2 aproveita uma informação que só existe no 2º turno: o erro que cada instituto cometeu três semanas antes, no 1º turno da mesma eleição. Medimos o erro na margem entre os dois primeiros e o dividimos por dois, para que ele fique na mesma escala de uma proporção de 2º turno:
\begin{equation}
e_j=\tfrac12\big[(\hat v^{\,1T}_{j,\text{PT}}-\hat v^{\,1T}_{j,\text{Adv}})-(v^{\,1T}_{\text{PT}}-v^{\,1T}_{\text{Adv}})\big],\qquad
w_j=\frac{1}{e_j^2+\sigma_0^2},\qquad \hat p_{\text{PT}}=\frac{\sum_j w_j\hat v_j}{\sum_j w_j}.
\end{equation}
O termo $\sigma_0$ impede que um instituto que acertou por sorte receba peso quase infinito. Institutos sem pesquisa de véspera no 1º turno recebem o erro mediano dos demais.

Validamos com \emph{expanding window}: 2014 é previsto com 2006 e 2010, 2018 com 2006 a 2014 e 2022 com 2006 a 2018, de modo que nenhuma previsão usa informação posterior. Um modelo mais complexo só substitui o mais simples se reduzir o MAE médio em mais de um erro-padrão da diferença, $\bar\Delta>s_\Delta/\sqrt3$. Em empate, fica o mais simples. Escolhemos primeiro o agregador e depois testamos cada ajuste isoladamente.

\begin{center}\small
\textbf{Tabela 1:} Backtest do 2º turno (MAE em p.p.; $\bar\Delta>0$ indica redução do erro).\\[1pt]
\begin{tabular}{lcccccl}
\toprule
\textbf{Etapa 1: agregador} & \textbf{2014} & \textbf{2018} & \textbf{2022} & \textbf{Média} & \textbf{$\bar\Delta$ / EP} & \textbf{Decisão} \\
\midrule
M0 média simples & 0,63 & 1,82 & 0,29 & 0,92 &  & referência \\
M1 recência e amostra ($h$=3) & 0,77 & 0,96 & 0,45 & 0,72 & +0,19 / 0,34 & não supera M0 \\
M2 peso pelo erro no 1º turno ($\sigma_0$=1,5) & 0,56 & 1,29 & 0,10 & 0,65 & +0,27 / 0,14 & \textbf{escolhido} \\
M2 peso pelo erro no 1º turno ($\sigma_0$=2,5) & 0,56 & 1,27 & 0,11 & 0,65 & +0,27 / 0,15 & supera, MAE maior \\
\midrule
\textbf{Etapa 2: ajustes sobre o M2} & & & & & & \\
\midrule
+ correção pelo erro do 1º turno ($\hat\rho$ MQO) & 0,29 & 1,40 & 0,01 & 0,57 & +0,08 / 0,11 & empate: fica a base \\
+ correção pelo erro do 1º turno ($\rho=0{,}2$) & 0,51 & 1,44 & 0,34 & 0,77 & $-$0,12 / 0,09 & piora \\
+ viés comum histórico do 2º turno ($k_\mu=3$) & 0,48 & 1,45 & 0,05 & 0,66 & $-$0,01 / 0,08 & empate: fica a base \\
+ 10\% do prior de transferência ($\tilde\tau$) & 1,61 & 1,62 & 0,12 & 1,12 & $-$0,47 / 0,31 & piora \\
\bottomrule
\end{tabular}
\\[1pt]
{\footnotesize O M2 erra menos que o M0 nas três eleições de teste e também em 2006 e 2010 (LOEO: \loeoVenc{} contra \loeoMZero). Fonte: elaboração própria com dados do TSE.}
\end{center}

O M2 com $\sigma_0=\sigmaZero$ reduziu o MAE de \maeMZero{} para \maeVenc{} ($\bar\Delta=\deltaVenc$, EP $=\seVenc$). Nenhum ajuste passou na regra. Corrigir cada pesquisa por uma fração do erro do seu instituto no 1º turno empatou com a base, e o viés comum histórico também: o erro de cada instituto no 1º turno não se repetiu de forma estável no 2º, possivelmente porque os institutos revisam suas ponderações entre os turnos. A previsão final é o M2 puro, arredondado a uma casa pelo método dos maiores restos, o que garante soma exata de 100,0.

\section{Uso do 1º turno}
O 1º turno entra de duas formas. A primeira é o peso de cada instituto no M2 (Figura 1b). Em 2026, todos os institutos com pesquisa de véspera superestimaram a vantagem de Lula, com erro mediano de \medErroUmT{} p.p. na metade da margem. A segunda é uma checagem independente pela transferência de votos. Definimos $\tau$ como a fração dos votos dos eliminados que o candidato do PT capturou entre os turnos:

\begin{center}\small
\textbf{Tabela 2:} Transferência aparente de votos do 1º para o 2º turno (\% dos válidos).\\[1pt]
\begin{tabular}{lcccccc}
\toprule
 & 2006 & 2010 & 2014 & 2018 & 2022 & \textbf{2026} \\
\midrule
PT no 1º turno & 48,6 & 46,9 & 41,6 & 29,3 & 48,4 & 45,2 \\
Demais candidatos no 1º turno & 9,8 & 20,5 & 24,9 & 24,7 & 8,4 & 7,8 \\
PT no 2º turno & 60,8 & 56,1 & 51,6 & 44,9 & 50,9 & ? \\
$\tau$ (fração capturada pelo PT) & 1,25 & 0,45 & 0,40 & 0,63 & 0,30 & $\tilde\tau=0,45$ \\
\bottomrule
\end{tabular}
\\[1pt]
{\footnotesize $\tau=(V^{2T}_{\text{PT}}-V^{1T}_{\text{PT}})/V^{1T}_{\text{demais}}$. Em 2006, $\tau>1$ porque Alckmin perdeu votos entre os turnos, e não houve só transferência. Fonte: elaboração própria com dados do TSE.}
\end{center}

Como 2006 distorce a média ($\bar\tau=\tauMedia$), usamos a mediana, $\tilde\tau=\tauMediana$. Aplicada ao 1º turno de 2026, ela leva Flávio a \priorMedianaFlavio\%, praticamente a previsão das pesquisas (\pFlavio\%). Lido ao contrário, nosso número supõe que Lula capture \tauImplLula\% dos \tercUmT{} p.p. dos eliminados e Flávio \tauImplFlavio\%, o padrão típico das eleições anteriores. Duas fontes independentes, as pesquisas de 2º turno e a transferência histórica sobre o 1º turno, chegam ao mesmo resultado. Não usamos a transferência como ajuste porque, misturada às pesquisas, ela piorou o backtest (Tabela 1): a dispersão de $\tau$ entre eleições é grande demais para servir de âncora.

\begin{figure}[h]
\centering
\begin{tikzpicture}
\begin{axis}[
  name=a, width=0.50\linewidth, height=4.8cm,
  xmin=0, xmax=20, ymin=49, ymax=56,
  xtick={0,5,10,15,20}, xticklabels={05/10,10/10,15/10,20/10,25/10},
  ylabel={Flávio, \% dos válidos}, label style={font=\scriptsize}, tick label style={font=\scriptsize},
  /pgf/number format/use comma, axis lines*=left, ymajorgrids, grid style={dashed,gray!30},
  legend style={font=\tiny, draw=none, fill=none, at={(0.5,1.02)}, anchor=south, legend columns=3, /tikz/every even column/.append style={column sep=4pt}}]
\addplot[only marks, mark=*, mark size=1.4pt, color=cinza!70] table[x=dia,y=flavio]{pesquisas_2026_2t.dat};
\addplot[thick, color=azul, domain=0:20] {\pFlavioDot};
\addplot[thick, dashed, color=vermelho, domain=0:20] {\flavioDoisCandDot};
\legend{Pesquisas de 2026, Previsão (\pFlavio), 1º turno entre os dois (\flavioDoisCand)}
\end{axis}
\begin{axis}[
  at={(a.east)}, anchor=west, xshift=2.3cm,
  width=0.42\linewidth, height=5.1cm, xbar, bar width=3.2pt, bar shift=0pt,
  xmin=0, ymin=-0.6, ymax=\nPesos.6,
  ytick={0,...,\nPesos}, yticklabels/.expand once=\rotulosPesos,
  xlabel={peso no M2 (\%)}, label style={font=\scriptsize}, tick label style={font=\tiny},
  /pgf/number format/use comma, axis lines*=left, xmajorgrids, grid style={dashed,gray!30},
  legend style={font=\tiny, draw=none, fill=none, at={(0.5,1.02)}, anchor=south, legend columns=2, /tikz/every even column/.append style={column sep=4pt}},
  legend image code/.code={\draw[#1] (0cm,-0.06cm) rectangle (0.22cm,0.08cm);}]
\addplot[fill=azul, draw=none] table[x=peso,y=y, restrict expr to domain={\thisrow{tem1t}}{1:1}]{pesos_2026.dat};
\addplot[fill=cinza!40, draw=none] table[x=peso,y=y, restrict expr to domain={\thisrow{tem1t}}{0:0}]{pesos_2026.dat};
\legend{com erro medido no 1º turno, erro mediano}
\end{axis}
\end{tikzpicture}
\caption{(a) Pesquisas de 2º turno em votos válidos, previsão final e a fatia de Flávio entre os dois primeiros no 1º turno. (b) Peso de cada instituto no M2: quem errou menos no 1º turno pesa mais. Fonte: elaboração própria.}
\end{figure}

\section{Abstenção, brancos e nulos}
No 2º turno já conhecemos esses valores no 1º turno da mesma eleição, e todos os métodos partem deles. Comparamos repetir o 1º turno com somar a variação entre turnos da eleição anterior ou a variação média, aplicar a razão média 2T/1T e decompor a abstenção por UF conforme haja 2º turno para governador. A validação e a regra são as mesmas da Tabela 1.

\begin{center}\small
\textbf{Tabela 3:} MAE médio (p.p.) em 2014, 2018 e 2022 e previsão.\\[1pt]
\begin{tabular}{lccc}
\toprule
\textbf{Método} & \textbf{Abstenção} & \textbf{Brancos} & \textbf{Nulos} \\
\midrule
Repetir o 1º turno & \textbf{1,01} & \textbf{0,93} & \textbf{0,93} \\
1º turno + variação da eleição anterior & 1,25 & 1,09 & 1,15 \\
1º turno + variação média & 1,67 & 1,01 & 1,11 \\
1º turno $\times$ razão média 2T/1T & 1,98 & 0,58 & 1,05 \\
Por UF, com e sem 2º turno estadual & 1,38 & n.a. & n.a. \\
\midrule
1º turno de 2026 (TSE) & 21,08 & 1,84 & 2,93 \\
\textbf{Previsão para o 2º turno} & \textbf{21,1} & \textbf{1,8} & \textbf{2,9} \\
\bottomrule
\end{tabular}
\\[1pt]
{\footnotesize Abstenção sobre o eleitorado apto; brancos e nulos sobre o comparecimento. Fonte: elaboração própria com dados do TSE.}
\end{center}

Nenhum método superou repetir o 1º turno por mais de um erro-padrão. O caso mais próximo é o dos brancos: eles caíram entre os turnos nas cinco eleições, e a razão média tem o menor MAE. Ainda assim ela errou mais que a persistência em 2018 e 2022, as duas eleições com nível de brancos próximo do atual, e por isso mantivemos o valor do 1º turno.
\textbf{Limitação.} Em 2026 só \govNVinteSeis{} UFs, com \govShareVinteSeis\% do eleitorado, têm 2º turno para governador, contra \govShareHist\% em média de 2006 a 2022. Sem campanha estadual, a abstenção pode subir mais que o histórico. O efeito medido foi pequeno (a abstenção subiu \deltaAbstCom{} p.p. nas UFs com disputa estadual e \deltaAbstSem{} nas demais), mas 2026 está fora da faixa observada.

\section{Limitações e reprodutibilidade}
Temos três eleições de teste, o que dá pouca potência à regra de decisão. Só \nInstUmT{} dos \nInst{} institutos de 2026 tinham pesquisa de véspera no 1º turno na nossa base, de modo que o M2 fica próximo da média simples (\sensMZero\%). Nas \nSens{} especificações alternativas que testamos, Flávio fica entre \sensMin\% e \sensMax\%. Essa faixa mostra quanto a previsão depende das escolhas de modelagem. Código, dados e testes: \url{https://github.com/pvaluche/trabalho-econometria-a1/tree/segundo-turno}.

\end{document}
